% Generated draft: requires mathematical and source review.
\documentclass[11pt]{article}
\usepackage[margin=25mm]{geometry}
\usepackage{fontspec}
\setmainfont{Latin Modern Roman}
\usepackage{amsmath,amssymb,mathtools,unicode-math}
\setmathfont{Latin Modern Math}
\usepackage{enumitem,needspace,xurl,hyperref}
\setlength{\parindent}{0pt}
\setlength{\parskip}{5pt}
\setlength{\emergencystretch}{3em}
\newcommand{\N}{\mathbb{N}}
\newcommand{\Z}{\mathbb{Z}}
\newcommand{\Q}{\mathbb{Q}}
\newcommand{\R}{\mathbb{R}}
\newcommand{\C}{\mathbb{C}}
\newcommand{\F}{\mathbb{F}}
\newcommand{\A}{\mathbb{A}}
\newcommand{\PP}{\mathbb{P}}
\newcommand{\E}{\mathbb{E}}
\newcommand{\eps}{\varepsilon}
\newcommand{\dd}{\,\mathrm{d}}
\begin{document}
\section*{UnsolvedMath ID: 30001729}
\textbf{Source problem code:} catalogue entry
\subsection{Definitions and mathematical statement}
Sexagesimal notation is positional notation with base \(60\). A Pythagorean triple is a triple \((a,b,c)\in\Z_{>0}^{3}\) satisfying \(a^{2}+b^{2}=c^{2}\); these numbers can represent the side lengths of a right triangle. These definitions explain terminology and do not establish how the tablet's entries were intended to be used.

The supplied source describes Plimpton 322 as an Old Babylonian tablet listing \(15\) sexagesimal rectangles or Pythagorean triples. Its historical-interpretive research question is:
\begin{quote}
What kind of mathematical table is Plimpton 322?
\end{quote}
The target is the mathematical classification or intended function of the table represented by these entries. The supplied record does not give the entries, column meanings, a definition of ``sexagesimal rectangles,'' or criteria for deciding among interpretations. Accordingly, this question cannot be converted from the supplied evidence into a uniquely specified theorem or classification problem. The description ``rectangles or Pythagorean triples'' is retained as source terminology, without asserting that the tablet explicitly lists complete integer triples or that either description establishes its historical purpose.
\subsection{Short English statement}
Determine what kind of mathematical table Plimpton 322 was intended to be. The source describes \(15\) sexagesimal rectangles or Pythagorean triples, but supplies insufficient data and interpretive criteria for a uniquely specified mathematical formulation.
\subsection{Sources}
\begin{itemize}
\item Local UnsolvedMath registry, supplied complete record: permanent ID \(30001729\), title ``Mathematical Classification of Plimpton 322,'' exact source problem code \texttt{catalogue entry}; collection: Oberwolfach Reports Open Problems.
\item Supplied report reference and literature source: \emph{Mini-Workshop: History of Numerical and Graphical Tables} (2011), Oberwolfach Reports, Report 12/2011, pp.\ 639--689, DOI \texttt{10.4171/owr/2011/12}. \url{https://doi.org/10.4171/owr/2011/12}. The supplied background describes Oberwolfach Reports as documenting mathematical workshops and their open-problem discussions. The record gives proposed year 2011 but identifies no individual proposer.
\item Supplied literature assessment, checked on \texttt{2026-08-21}: status ``open''; the question is characterized as a specific historical-interpretive research question for which the located literature does not establish a unique settled answer. This is supplied dated curation, not independently verified. The listed source checked is Oberwolfach Report 12/2011, linked above. The supplied background requests independent verification of time-sensitive status and cited claims.
\item The registry labels the statement \texttt{corrected\_verified} and reports that the official report confirms the historiographical research question and that an extraction label was removed. These are supplied verification assertions, not independent verification performed for this draft.
\end{itemize}
\end{document}
